\section{Related Literature}

The paper engages three strands of literature.

\subsection{Technology Adoption and Productivity}

The first strand studies how broad technologies shape productivity and organizational performance. Classical work emphasizes that returns depend on complementary assets such as managerial quality, human capital, and organizational redesign \cite{brynjolfssonHitt2000,griliches1998}. Under this view, technology does not mechanically raise output; realized gains depend on diffusion frictions and institutional capacity \cite{bresnahanTrajtenberg1995,cominMestieri2018}. These insights are directly relevant for AI, where implementation quality and complementary capabilities vary markedly across countries.

\subsection{AI, Automation, and Aggregate Outcomes}

The second strand focuses on AI and automation. Conceptual contributions suggest that AI may generate sizable long-run growth effects, but with transitional frictions in labor reallocation and capital adjustment \cite{acemogluRestrepo2018,brynjolfssonRockSyverson2019,aghionJonesJones2017,acemogluAutor2011}. Policy-oriented analyses similarly highlight both opportunity and adjustment cost, especially where digital infrastructure and skills are unevenly distributed \cite{imf2024,oecd2023ai,korinekStiglitz2021}. Macro evidence remains mixed, partly because AI intensity is measured imperfectly and partly because diffusion lags can mute short-run growth responses.

Recent evidence from 2023--2024 reinforces this caution: the policy and measurement literature increasingly documents rapid adoption alongside uneven capability readiness, implying that cross-country aggregate elasticities are likely to be moderate and context-dependent rather than uniform \cite{oecd2023ai,imf2024}. Micro- and task-level studies also show meaningful productivity gains from generative AI tools, but those estimates are typically identified in narrower settings with shorter horizons \cite{noyZhang2023,brynjolfssonLiRaymond2023,eloundouEtAl2023}. This also clarifies why the academic macro literature can appear sparse at this stage: available AI measures and sufficiently long post-adoption panels are still catching up to the speed of diffusion.

\subsection{Panel Identification and Inference}

The third strand concerns empirical strategy in panel settings. Fixed-effects estimators are widely used to absorb time-invariant heterogeneity, but they do not by themselves solve omitted time-varying confounding. Two-way fixed effects, clustering choices, and specification diagnostics affect both point estimates and inference quality \cite{wooldridge2010,wooldridge2021}. Recent practice therefore emphasizes robustness portfolios rather than reliance on a single preferred model.

\subsection{Positioning of the Present Study}

Relative to this literature, the main contribution here is methodological transparency with disciplined interpretation. The paper links data checks, model estimation, and manuscript outputs in a reproducible workflow, and then evaluates whether conclusions persist across pre-specified robustness and sensitivity models. The empirical objective is not to claim final causality, but to identify stable conditional associations that can guide future designs with stronger exogenous variation.
